\chapter{Clauze, Rezoluție și Prolog}

\section{Continuare FNC și FND}

Exercițiile de la forme normale sunt foarte algoritmice și există foarte mare șansă să aveți la examen acest tip de exercițiu. Așa că e bine să vă obișnuiți să lucrați cu ele. Există două tipuri de exerciții: transformarea unei formule în FNC/FND prin transformare sintactică (adică al doilea algoritm prezentat în capitolul anterior) și prin aducerea la funcție booleană (adică compunerea tabelului de adevăr a formululei și aplicarea primului algoritm din capitolul anterior).

\begin{ex}
Să se aducă următoarele formule la cele două forme normale prin transformări sintactice:
\begin{enumerate}
    \item[(i)] $((v_0 \rightarrow v_1) \land v_1) \rightarrow v_0$;
    \item[(ii)] $(v_1 \lor \neg v_4) \rightarrow (\neg v_2 \rightarrow v_3)$.
\end{enumerate}
\end{ex}

\begin{solutie}

Vom folosi algoritmul detaliat în \ref{alg2:fn}.

\vspace{0.3cm}
\noindent \textbf{Pentru formula (i):} $\alpha = ((v_0 \rightarrow v_1) \land v_1) \rightarrow v_0$

\begin{itemize}
    \item \textit{Pasul 1 (Eliminarea implicațiilor):}
          Aplicăm regula $\varphi \rightarrow \psi \sim \neg \varphi \lor \psi$:
          \[ \alpha \sim \neg ((\neg v_0 \lor v_1) \land v_1) \lor v_0 \]
    \item \textit{Pasul 2 (Aplicarea legilor De Morgan și absorbția):}
          Introducem negația în paranteză folosind De Morgan:
          \[ \alpha \sim (\neg (\neg v_0 \lor v_1) \lor \neg v_1) \lor v_0 \]
          Aplicăm din nou De Morgan pentru termenul $\neg (\neg v_0 \lor v_1)$:
          \[ \alpha \sim ((v_0 \land \neg v_1) \lor \neg v_1) \lor v_0 \]
    \item \textit{Pasul 3 (Aplicarea distributivității și a idempotenței):}
	      Aplicăm distributivitatea disjuncției față de conjuncție pentru subformula $((v_0 \land \neg v_1) \lor \neg v_1)$:
	      \[ \alpha \sim ((v_0 \lor \neg v_1) \land (\neg v_1 \lor \neg v_1)) \lor v_0 \]
	      Aplicăm legea idempotenței pentru disjuncție ($\neg v_1 \lor \neg v_1 \sim \neg v_1$):
	      \[ \alpha \sim ((v_0 \lor \neg v_1) \land \neg v_1) \lor v_0 \]

	\item \textit{Pasul 4 (A doua aplicare a distributivității):}
	      Pentru a finaliza aducerea la forma normală, distribuim disjuncția $(\lor v_0)$ peste conjuncția rămasă:
	      \[ \alpha \sim ((v_0 \lor \neg v_1) \lor v_0) \land (\neg v_1 \lor v_0) \]
	      Folosind comutativitatea, asociativitatea și idempotența ($v_0 \lor v_0 \sim v_0$) în prima paranteză, obținem:
	      \[ \alpha \sim (v_0 \lor \neg v_1) \land (v_0 \lor \neg v_1) \]
	      În final, aplicăm idempotența pentru întreaga expresie ($A \land A \sim A$), rezultând o formulă care e atât în FNC, cât și
          în FND:
	      \[ \alpha \sim v_0 \lor \neg v_1 \]
\end{itemize}


\noindent \textbf{Pentru formula (ii):} $\beta = (v_1 \lor \neg v_4) \rightarrow (\neg v_2 \rightarrow v_3)$

\begin{itemize}
    \item \textit{Pasul 1 (Eliminarea implicațiilor):}
          Mai întâi eliminăm implicația din paranteza dreaptă:
          \[ \beta \sim (v_1 \lor \neg v_4) \rightarrow (\neg\neg v_2 \lor v_3) \sim (v_1 \lor \neg v_4) \rightarrow (v_2 \lor v_3) \]
          Apoi eliminăm implicația principală:
          \[ \beta \sim \neg (v_1 \lor \neg v_4) \lor (v_2 \lor v_3) \]
    \item \textit{Pasul 2 (Aplicarea legilor De Morgan):}
          \[ \beta \sim (\neg v_1 \land \neg\neg v_4) \lor v_2 \lor v_3 \sim (\neg v_1 \land v_4) \lor v_2 \lor v_3 \]
          În acest stadiu, formula este deja o disjuncție de conjuncții, așadar avem forma FND.
    \item \textit{Pasul 3 (Aducerea la FNC prin distributivitate):}
          Pentru a obține FNC, distribuim disjuncția $(\lor v_2 \lor v_3)$ peste conjuncția $(\neg v_1 \land v_4)$:
          \[ \beta \sim (\neg v_1 \lor v_2 \lor v_3) \land (v_4 \lor v_2 \lor v_3) \]
\end{itemize}


\end{solutie}


\begin{ex}
\textbf Să se aducă formula $\varphi = (v_0 \rightarrow v_1) \rightarrow v_2$ la cele două forme normale trecându-se prin funcția booleană asociată (i.e. metoda tabelului).
\end{ex}

\begin{solutie}
Construim tabelul de adevăr pentru funcția booleană asociată formulei $\varphi$. Variabilele sunt $v_0, v_1, v_2$.

\vspace{0.3cm}

\begin{minipage}[c]{0.35\textwidth}
	\renewcommand{\arraystretch}{1.2}
	\begin{tabular}{c|c|c|c|c}
		$v_0$ & $v_1$ & $v_2$ & $v_0 \to v_1$ & $\varphi$ \\
		\hline
		0 & 0 & 0 & 1 & \textcolor{blue}{0} \\
		0 & 0 & 1 & 1 & \textcolor{magenta}{1} \\
		0 & 1 & 0 & 1 & \textcolor{blue}{0} \\
		0 & 1 & 1 & 1 & \textcolor{magenta}{1} \\
		1 & 0 & 0 & 0 & \textcolor{magenta}{1} \\
		1 & 0 & 1 & 0 & \textcolor{magenta}{1} \\
		1 & 1 & 0 & 1 & \textcolor{blue}{0} \\
		1 & 1 & 1 & 1 & \textcolor{magenta}{1} \\
	\end{tabular}
\end{minipage}
\begin{minipage}[c]{0.6\textwidth}
	\begin{align*}
		\textcolor{blue}{D_1}    & = \textcolor{blue}{v_0 \lor v_1 \lor v_2} \\
		\textcolor{magenta}{C_1} & = \textcolor{magenta}{\neg v_0 \land \neg v_1 \land v_2} \\
		\textcolor{blue}{D_2}    & = \textcolor{blue}{v_0 \lor \neg v_1 \lor v_2} \\
		\textcolor{magenta}{C_2} & = \textcolor{magenta}{\neg v_0 \land v_1 \land v_2} \\
		\textcolor{magenta}{C_3} & = \textcolor{magenta}{v_0 \land \neg v_1 \land \neg v_2} \\
		\textcolor{magenta}{C_4} & = \textcolor{magenta}{v_0 \land \neg v_1 \land v_2} \\
		\textcolor{blue}{D_3}    & = \textcolor{blue}{\neg v_0 \lor \neg v_1 \lor v_2} \\
		\textcolor{magenta}{C_5} & = \textcolor{magenta}{v_0 \land v_1 \land v_2}
	\end{align*}
\end{minipage}

\vspace{0.5cm}

\noindent Extragem \textbf{Forma Normală Disjunctivă (FND)} prin disjuncția conjuncțiilor $C_i$ corespunzătoare liniilor unde $\varphi = \textcolor{magenta}{1}$:

\vspace{0.2cm}
\noindent $\varphi^{FND} = C_1 \lor C_2 \lor C_3 \lor C_4 \lor C_5$ \\
$\varphi^{FND} = (\neg v_0 \land \neg v_1 \land v_2) \lor (\neg v_0 \land v_1 \land v_2) \lor (v_0 \land \neg v_1 \land \neg v_2) \lor (v_0 \land \neg v_1 \land v_2) \lor (v_0 \land v_1 \land v_2)$

\vspace{0.5cm}

\noindent Extragem \textbf{Forma Normală Conjunctivă (FNC)} prin conjuncția disjuncțiilor $D_i$ corespunzătoare liniilor unde $\varphi = \textcolor{blue}{0}$ (se neagă valoarea de adevăr a variabilelor de pe linie):


\vspace{0.2cm}
\noindent $\varphi^{FNC} = D_1 \land D_2 \land D_3$ \\
$\varphi^{FNC} = (v_0 \lor v_1 \lor v_2) \land (v_0 \lor \neg v_1 \lor v_2) \land (\neg v_0 \lor \neg v_1 \lor v_2)$

\end{solutie}

\section{Clauze și forma clauzală}

Până acum am lucrat cu formule în FNC/FND. Pentru algoritmi de satisfiabilitate și pentru rezoluție, este util să tratăm formulele în FNC drept \textbf{mulțimi de clauze}. Asta simplifică mult calculele, pentru că ordinea literalilor nu contează, iar duplicatele se pot elimina imediat.

\subsection{Definiții de bază}

\begin{definition}
	O \textbf{clauză} este o mulțime finită de literali:
	\[
		C = \{L_1, L_2, \dots, L_n\}.
	\]

	\noindent Dacă \(n = 0\), obținem \textbf{clauza vidă}, notată \( \square \), adică \( \square = \emptyset \).
\end{definition}

Dacă \(L\) este un literal, notăm literalul complementar cu \(L^c\):
\[
	L^c =
	\begin{cases}
		\neg v, & \text{dacă } L = v \\
		v, & \text{dacă } L = \neg v
	\end{cases}
\]

Intuitiv, o clauză nevidă se citește ca o disjuncție de literali. De exemplu:
\[
	\{v_1, \neg v_3, v_5\} \quad \leftrightarrow \quad v_1 \lor \neg v_3 \lor v_5.
\]

\begin{definition}
	Fie \(C\) o clauză și \(e: V \to \{0,1\}\) o evaluare.
	Spunem că \(e\) satisface clauza \(C\), și scriem \(e \models C\), dacă există un literal \(L \in C\) cu \(e \models L\).

	\noindent Pentru o mulțime finită de clauze \(S = \{C_1, \dots, C_m\}\), scriem \(e \models S\) dacă \(e \models C_i\) pentru orice \(i\). Deci \(S\) se citește ca o conjuncție de clauze.
\end{definition}

\begin{definition}
	Fie \(C\) o clauză și \(S\) o mulțime finită de clauze.
	\begin{enumerate}
		\item Spunem că \(C\) este \textbf{satisfiabilă} dacă există o evaluare \(e\) cu \(e \models C\). Spunem că \(C\) este \textbf{validă} dacă pentru orice evaluare \(e\), avem \(e \models C\).
		\item Spunem că \(S\) este \textbf{satisfiabilă} dacă există o evaluare \(e\) cu \(e \models S\). Spunem că \(S\) este \textbf{validă} dacă pentru orice evaluare \(e\), avem \(e \models S\).
	\end{enumerate}
\end{definition}

\begin{prop}
    \noindent
	\begin{enumerate}
		\item Orice clauză nevidă este satisfiabilă.
		\item Clauza vidă \( \square \) este nesatisfiabilă.
		\item Dacă o clauză conține simultan \(L\) și \(L^c\), atunci este validă (trivială).
	\end{enumerate}
\end{prop}

\subsection{Legătura dintre clauze și FNC}

Pentru o formulă în FNC
\[
	\varphi = \bigwedge_{i=1}^{n}\left(\bigvee_{j=1}^{k_i} L_{i,j}\right),
\]
asociem forma clauzală
\[
	S_\varphi = \{C_1, \dots, C_n\}, \quad C_i = \{L_{i,1}, \dots, L_{i,k_i}\},
\]
unde într-o clauză păstrăm fiecare literal o singură dată.

\begin{ex}
	Fie
	\[
		\varphi = (v_1 \lor \neg v_2 \lor v_1)\land (\neg v_1 \lor v_3)\land (v_2).
	\]
	Să se scrie forma clauzală a formulei.
\end{ex}

\begin{solutie}
	Transformăm fiecare paranteză într-o clauză și eliminăm duplicatele:
	\[
		S_\varphi = \{\{v_1, \neg v_2\}, \{\neg v_1, v_3\}, \{v_2\}\}.
	\]

	Observația importantă este că \( \varphi \) și \( S_\varphi \) sunt semantic echivalente: au exact aceleași modele.
\end{solutie}

\subsection{Conversia inversă: din clauze la formula $\varphi_S$}

Conversia se face în doi pași: mai întâi transformăm fiecare clauză într-o disjuncție, apoi facem conjuncția tuturor acestor formule.

\begin{definition}
	Pentru o clauză \(C\), formula asociată \( \varphi_C \) este:
	\begin{enumerate}
		\item dacă \(C = \{L_1, \dots, L_n\}\), cu \(n \geq 1\), atunci
		      \[
			      \varphi_C := L_1 \lor \dots \lor L_n;
		      \]
		\item dacă \(C = \square\), atunci putem lua
		      \[
			      \varphi_{\square} := v_0 \land \neg v_0.
		      \]
	\end{enumerate}
\end{definition}

\begin{definition}
	Pentru o mulțime finită de clauze \(S\), formula asociată \( \varphi_S \) este:
	\begin{enumerate}
		\item dacă \(S = \{C_1, \dots, C_m\}\), cu \(m \geq 1\), atunci
		      \[
			      \varphi_S := \bigwedge_{i=1}^{m} \varphi_{C_i};
		      \]
		\item dacă \(S = \emptyset\), atunci putem lua
		      \[
			      \varphi_{\emptyset} := v_0 \lor \neg v_0.
		      \]
	\end{enumerate}
\end{definition}

\begin{note}
	Formula \( \varphi_S \) nu este unică sintactic (ordinea literalilor/clauzelor poate diferi), dar toate variantele sunt echivalente semantic.
\end{note}

\begin{prop}
	Pentru orice evaluare \(e\), avem:
	\[
		e \models S \quad \text{ddacă} \quad e \models \varphi_S.
	\]
\end{prop}

\begin{ex}
	Fie
	\[
		S = \{\{v_1, \neg v_2\}, \{v_3\}, \{\neg v_1, v_4\}\}.
	\]
	Scrieți formula \( \varphi_S \).
\end{ex}

\begin{solutie}
	Asociem formulele pe clauze:
	\[
		\varphi_{C_1} = v_1 \lor \neg v_2, \quad \varphi_{C_2} = v_3, \quad \varphi_{C_3} = \neg v_1 \lor v_4.
	\]
	Deci:
	\[
		\varphi_S = (v_1 \lor \neg v_2)\land v_3 \land (\neg v_1 \lor v_4).
	\]
\end{solutie}

\section{Rezoluția pe clauze}

Rezoluția este o regulă de inferență specială pentru clauze. Ea este folosită ca să verificăm dacă un set de condiții este compatibil sau nu. Pe scurt, dacă două clauze se „bat cap în cap” pe aceeași variabilă (una conține \(L\), cealaltă conține \(L^c\)), eliminăm perechea contradictorie și păstrăm ce rămâne.

Exemplu:
\[
	\{p, q\} \ \text{și}\ \{\neg p, r\} \ \Longrightarrow \ \{q, r\}.
\]
Am eliminat \(p\) și \(\neg p\), pentru că sunt complementare.

Dacă, după mai mulți astfel de pași, ajungem la clauza vidă \( \square \), înseamnă că am obținut contradicție, deci clauzele inițiale nu pot fi satisfăcute împreună. Exact de aceea rezoluția este utilă în verificare automată, în solvere SAT și în mecanismul logic din PROLOG.

\begin{definition}
	Fie \(C_1, C_2\) două clauze. O clauză \(R\) se numește \textbf{rezolvent} al lor dacă există un literal \(L\) astfel încât
	\[
		L \in C_1, \quad L^c \in C_2, \quad R = (C_1 \setminus \{L\}) \cup (C_2 \setminus \{L^c\}).
	\]
\end{definition}

\begin{ex}
	Fie \(C_1 = \{v_1, \neg v_2, v_4\}\) și \(C_2 = \{v_2, v_3\}\).
	Să se calculeze un rezolvent.
\end{ex}

\begin{solutie}
	Alegem literalul \(L = \neg v_2\), deci \(L^c = v_2\).
	\[
		R = (C_1 \setminus \{\neg v_2\}) \cup (C_2 \setminus \{v_2\}) = \{v_1, v_4, v_3\}.
	\]
	Acesta este un rezolvent al lui \(C_1\) și \(C_2\).
\end{solutie}

\begin{prop}[Corectitudinea rezoluției]
	Dacă dintr-o mulțime de clauze \(S\) se poate deriva clauza vidă \( \square \) prin pași de rezoluție, atunci \(S\) este nesatisfiabilă.
\end{prop}

\begin{ex}
	Arătați prin rezoluție că
	\[
		S = \{\{v_1\}, \{\neg v_1, v_2\}, \{\neg v_2\}\}
	\]
	este nesatisfiabilă.
\end{ex}

\begin{solutie}
	Notăm
	\[
		C_1=\{v_1\},\quad C_2=\{\neg v_1, v_2\},\quad C_3=\{\neg v_2\}.
	\]
	Rezolvăm \(C_1\) cu \(C_2\) după literalul \(v_1\):
	\[
		C_4 = \{v_2\}.
	\]
	Rezolvăm \(C_4\) cu \(C_3\) după literalul \(v_2\):
	\[
		C_5 = \square.
	\]
	Am derivat clauza vidă, deci \(S\) este nesatisfiabilă.
\end{solutie}

\section{Algoritmul Davis--Putnam (DP)}

Algoritmul Davis--Putnam decide satisfiabilitatea unei mulțimi finite de clauze prin \textbf{eliminare de variabile} folosind rezolvenți.

\subsection*{Algoritmul DP (scriere standard)}

\noindent \textbf{Intrare:} \(S\) mulțime finită nevidă de clauze netriviale.
\[
	i := 1, \quad S_1 := S.
\]

\noindent \textit{Pi.1} Fie \(x_i\) o variabilă care apare în \(S_i\). Definim:
\[
	T_i^1 := \{C \in S_i \mid x_i \in C\}, \qquad
	T_i^0 := \{C \in S_i \mid \neg x_i \in C\}.
\]

\noindent \textit{Pi.2} Dacă \(T_i^1 \neq \emptyset\) și \(T_i^0 \neq \emptyset\), atunci
\[
	U_i := \{(C_1 \setminus \{x_i\}) \cup (C_0 \setminus \{\neg x_i\}) \mid C_1 \in T_i^1,\ C_0 \in T_i^0\}.
\]
altfel \(U_i := \emptyset\).

\noindent \textit{Pi.3} Definim:
\[
	S_{i+1}' := (S_i \setminus (T_i^0 \cup T_i^1)) \cup U_i,
\]
\[
	S_{i+1} := S_{i+1}' \setminus \{C \in S_{i+1}' \mid C \text{ trivială}\}.
\]

\noindent \textit{Pi.4} Facem testul de oprire:
\[
	\text{dacă } S_{i+1} = \emptyset \text{ atunci } S \text{ este satisfiabilă;}
\]
\[
	\text{altfel dacă } \square \in S_{i+1} \text{ atunci } S \text{ este nesatisfiabilă;}
\]
\[
	\text{altfel } i := i+1 \text{ și revenim la Pi.1}.
\]

\subsection*{Explicația pașilor}

Intuitiv, la fiecare iterație alegem o variabilă \(x_i\) și „o eliminăm” din sistem, păstrând doar informația relevantă pentru satisfiabilitate.

\noindent \textbf{Pasul Pi.1.} Împărțim clauzele din \(S_i\) în două grupuri: cele care conțin \(x_i\) (acestea formează \(T_i^1\)) și cele care conțin \(\neg x_i\) (acestea formează \(T_i^0\)). Clauzele care nu conțin nici \(x_i\), nici \(\neg x_i\), rămân pe loc.

\noindent \textbf{Pasul Pi.2.} Dacă avem clauze în ambele grupuri (\(T_i^1\) și \(T_i^0\)), construim toate rezolventele posibile între ele. Mulțimea acestor rezolvente este \(U_i\). Dacă unul dintre grupuri este vid, nu avem ce rezolva și luăm \(U_i=\emptyset\).

\noindent \textbf{Pasul Pi.3.} Eliminăm din \(S_i\) toate clauzele care conțin \(x_i\) sau \(\neg x_i\), apoi adăugăm rezolventele din \(U_i\). După aceea, eliminăm clauzele triviale (cele care conțin simultan un literal și complementul său), deoarece nu ajută la decizie.

\noindent \textbf{Pasul Pi.4.} Dacă după această simplificare obținem mulțimea vidă de clauze, problema este satisfiabilă. Dacă apare clauza vidă \( \square \), problema este nesatisfiabilă. În rest, repetăm aceeași procedură cu o nouă variabilă.

\begin{ex}
	Fie
	\[
		S = \{\{v_1, \neg v_3\}, \{v_2, v_1\}, \{v_2, \neg v_1, v_3\}\}.
	\]
	Să se aplice DP.
\end{ex}

\begin{solutie}
	\noindent \(S = \{\{v_1, \neg v_3\}, \{v_2, v_1\}, \{v_2, \neg v_1, v_3\}\}\), \(i := 1\), \(S_1 := S\).

	\noindent \textbf{P1.1} \(x_1 := v_3\), \(T_1^1 := \{\{v_2, \neg v_1, v_3\}\}\), \(T_1^0 := \{\{v_1, \neg v_3\}\}\).\\
	\textbf{P1.2} \(U_1 := \{\{v_2, \neg v_1, v_1\}\}\).\\
	\textbf{P1.3} \(S_2' := \{\{v_2, v_1\}, \{v_2, \neg v_1, v_1\}\}\), \(S_2 := \{\{v_2, v_1\}\}\).\\
	\textbf{P1.4} \(i := 2\) and go to P2.1.

	\vspace{0.2cm}
	\noindent \textbf{P2.1} \(x_2 := v_2\), \(T_2^1 := \{\{v_2, v_1\}\}\), \(T_2^0 := \emptyset\).\\
	\textbf{P2.2} \(U_2 := \emptyset\).\\
	\textbf{P2.3} \(S_3 := \emptyset\).\\
	\textbf{P2.4} \(S\) este satisfiabilă.
\end{solutie}

\begin{ex}
	Fie
	\[
		S = \{\{\neg v_1, v_2, \neg v_4\}, \{\neg v_3, \neg v_2\}, \{v_1, v_3\}, \{v_1\}, \{v_3\}, \{v_4\}\}.
	\]
	Să se aplice DP și să se decidă satisfiabilitatea.
\end{ex}

\begin{solutie}
	\noindent \(S = \{\{\neg v_1, v_2, \neg v_4\}, \{\neg v_3, \neg v_2\}, \{v_1, v_3\}, \{v_1\}, \{v_3\}, \{v_4\}\}\),\\
	\noindent \(i := 1\), \(S_1 := S\).

	\noindent \textbf{P1.1} \(x_1 := v_1\), \(T_1^1 := \{\{v_1, v_3\}, \{v_1\}\}\), \(T_1^0 := \{\{\neg v_1, v_2, \neg v_4\}\}\).\\
	\textbf{P1.2} \(U_1 := \{\{v_3, v_2, \neg v_4\}, \{v_2, \neg v_4\}\}\).\\
	\textbf{P1.3} \(S_2 := \{\{\neg v_3, \neg v_2\}, \{v_3\}, \{v_4\}, \{v_3, v_2, \neg v_4\}, \{v_2, \neg v_4\}\}\).\\
	\textbf{P1.4} \(i := 2\) and go to P2.1.

	\vspace{0.2cm}
	\noindent \textbf{P2.1} \(x_2 := v_2\), \(T_2^1 := \{\{v_3, v_2, \neg v_4\}, \{v_2, \neg v_4\}\}\),\\
	\noindent \(T_2^0 := \{\{\neg v_3, \neg v_2\}\}\).\\
	\textbf{P2.2} \(U_2 := \{\{v_3, \neg v_4, \neg v_3\}, \{\neg v_4, \neg v_3\}\}\).\\
	\textbf{P2.3} \(S_3 := \{\{v_3\}, \{v_4\}, \{\neg v_4, \neg v_3\}\}\).\\
	\textbf{P2.4} \(i := 3\) and go to P3.1.

	\vspace{0.2cm}
	\noindent \textbf{P3.1} \(x_3 := v_3\), \(T_3^1 := \{\{v_3\}\}\), \(T_3^0 := \{\{\neg v_4, \neg v_3\}\}\).\\
	\textbf{P3.2} \(U_3 := \{\{\neg v_4\}\}\).\\
	\textbf{P3.3} \(S_4 := \{\{v_4\}, \{\neg v_4\}\}\).\\
	\textbf{P3.4} \(i := 4\) and go to P4.1.

	\vspace{0.2cm}
	\noindent \textbf{P4.1} \(x_4 := v_4\), \(T_4^1 := \{\{v_4\}\}\), \(T_4^0 := \{\{\neg v_4\}\}\).\\
	\textbf{P4.2} \(U_4 := \{\square\}\).\\
	\textbf{P4.3} \(S_5 := \{\square\}\).\\
	\textbf{P4.4} \(S\) este nesatisfiabilă.
\end{solutie}

\newpage

\section{Prolog findall/setof/bagof}

In prolog, in mod normal, cand pui un query, sistemul iti returneaza prima solutie
si continua cautarea pentru mai multe solutii doar daca vrem noi.

Predicatele \texttt{findall}, \texttt{bagof}, \texttt{setof} sunt folosite atunci cand
vrem sa colectam \textbf{toate} solutiile intr-o lista.

Diferentele intre ele stau in modul in care trateaza duplicatele, sortarea, lipsa solutiilor si,
cel mai important, variabilele libere.

Vom analiza cum se comporta fiecare incepand cu urmatoarele fapte simple:

\begin{lstlisting}
angajat(dana, it).
angajat(alex, it).
angajat(maria, hr).
angajat(alex, it). % Duplicat intentionat
\end{lstlisting}

\subsection*{findall(sablon, scop, listaRezultat)}
Este cel mai simplu si intuitiv si cel mai des folosit in practica. Face exact ce spune:
gaseste toate solutiile si le pune intr-o lista.

Pastreaza duplicate, nu sorteaza si \textbf{ignora variabilele libere}.

Cand nu gaseste nimic returneaza o lista goala (nu da false).

\textbf{Exemple:}

\texttt{?- findall(Nume, angajat(Nume, pr), R).}

\texttt{R = [].}

\vspace{\baselineskip}

\texttt{?- findall(Nume, angajat(Nume, Dept), R).}

\texttt{R = [dana, alex, maria, alex].}

\textit{Se observa cum a ignorat complet variabila libera Dept, a pastrat duplicatul 'alex' si nu a sortat nimic.}


\subsection*{bagof(sablon, scop, listaRezultat)}
Predicatul \texttt{bagof} este putin mai inteligent in sensul ca \textbf{grupeaza} rezultatele in functie de variabilele
libere (variabilele care sunt in \texttt{scop} dar nu sunt in \texttt{sablon}).

Pastreaza duplicate, nu sorteaza, si daca exista variabile libere, Prolog \textbf{va returna mai multe liste},
cate una pentru fiecare valoare posibila a variabilei libere.

Cand nu gaseste nimic \textbf{esueaza} (returneaza false).

\textbf{Exemple:}

\texttt{?- bagof(Nume, angajat(Nume, pr), R).}

\texttt{false.}

\vspace{\baselineskip}

\texttt{?- bagof(Nume, angajat(Nume, Dept), R).}

\texttt{Dept = hr,}

\texttt{R = [maria] ;}

\texttt{Dept = it,}

\texttt{R = [dana, alex, alex].}

\newpage

De aici putem sa incepem sa si combinam, de exemplu daca vrem o lista de tupluri de tipul:

\texttt{[(numeDepartament, [angajati in departament])]}

Putem face:

\texttt{?- findall((Dept, L), bagof(Nume, angajat(Nume, Dept), L), R).}

\texttt{R = [(hr, [maria]), (it, [dana, alex, alex])].}



\subsection*{setof(sablon, scop, listaRezultat)}
Predicatul \texttt{setof} este cel mai strict. El modeleaza o multime matematica (un set), si, in general, elementele unei multimi sunt ordonate.

\textbf{Elimina} duplicate, \textbf{sorteaza} elementele (alfabetic sau numeric), si similar cu bagof,
grupeaza dupa variabilele libere si esueaza cand nu gaseste nimic

\textbf{Exemple:}

\texttt{?- setof(Nume, angajat(Nume, pr), R).}

\texttt{false.}

\vspace{\baselineskip}

\texttt{?- setof(Nume, angajat(Nume, Dept), R).}

\texttt{Dept = hr,}

\texttt{R = [maria] ;}

\texttt{Dept = it,}

\texttt{R = [alex, dana].}

\textit{Se observa cum a eliminat duplicatul 'alex' si a sortat listele.}

\vspace{\baselineskip}

\texttt{?- findall((Dept, L), setof(Nume, angajat(Nume, Dept), L), R).}

\texttt{R = [(hr, [maria]), (it, [alex, dana])].}




\subsection*{Operatorul \^{}}

Operatorul \texttt{\^{}} este folosit in \texttt{setof} si \texttt{bagof} sa ignore o variabila libera data (similar cu \texttt{findall}
dar putem controla ce variabile libere sunt ignorate).

\textbf{Exemple:}

\texttt{?- setof(Nume, Dept\^{}angajat(Nume, Dept), R).}

\texttt{R = [alex, dana, maria].}

In cazul in care avem mai multe variabile libere pe care vrem sa le ignoram putem face:

\texttt{?- setof(Nume, Firma\^{}Dept\^{}lucreaza(Nume, Oras, Firma, Dept), R).}





\subsection*{Exercitii rezolvate}

\subsubsection*{Exercitiul 1}

Avand ca si context urmatoarele fapte:

\begin{lstlisting}
% inventar(Jucator, Obiect, ValoareAur).
inventar(arthas, sabie, 150).
inventar(arthas, potiune_viata, 10).
inventar(arthas, scut, 50).
inventar(jaina, toiag, 200).
inventar(jaina, potiune_mana, 15).
inventar(jaina, potiune_mana, 15). % Are doua potiuni
\end{lstlisting}

Definiti predicatul \texttt{avereJucator/2} care primeste un jucator ca prim parametru
si in al doilea parametru se afla aurul total pe care il are acel jucator.

Rulati si confirmati:

\texttt{?- avere\_jucator(arthas, R).}

\texttt{R = 210.}

\texttt{?- avere\_jucator(jaina, R).}

\texttt{R = 230.}

\textbf{Atentie!} Ce credeti ca se intampla daca rulam:

\texttt{?- avere\_jucator(X, 210).}

\texttt{?- avere\_jucator(X, 440).}

\texttt{?- avere\_jucator(X, R).}


\subsubsection*{Solutii:}

\subsubsection*{Solutia 1 findall:}

\begin{lstlisting}
avere_jucator(J, R) :-
    findall(V, inventar(J, _, V), L),
    suma_lista(L, R).
\end{lstlisting}

In aceasta solutie daca rulam:

\texttt{?- avere\_jucator(X, 210).}

\texttt{false.}

\texttt{?- avere\_jucator(X, 440).}

\texttt{true.}

\texttt{?- avere\_jucator(X, R).}

\texttt{R = 440.}

De ce? Pentru ca, daca ii dam sa caute jucatorul, atunci J o sa fie o variabila libera. In \texttt{findall}, variabilele libere sunt ignorate.
Prin urmare, linia cu \texttt{findall(V, inventar(J, \_, V), L),} va returna mereu lista cu toate valorile: \texttt{[150, 10, 50, 200, 15, 15]} si suma va fi 440.


\subsubsection*{Solutia 2 bagof:}
\begin{lstlisting}
avere_jucator(J, R) :-
    bagof(V, Obj^inventar(J, Obj, V), L),
    suma_lista(L, R).
\end{lstlisting}

In aceasta solutie daca rulam:

\texttt{?- avere\_jucator(X, 210).}

\texttt{X = arthas.}

\texttt{?- avere\_jucator(X, 440).}

\texttt{false.}

\texttt{?- avere\_jucator(X, R).}

\texttt{X = arthas,}

\texttt{R = 210 ;}

\texttt{X = jaina,}

\texttt{R = 230.}

\newpage

De ce? Pentru ca \texttt{bagof} va grupa dupa J cand este variabila libera si
cand verifica va lua toate gruparile posibile si forma listele:
\texttt{(arthas, [150, 10, 50])} si \texttt{(jaina, [200, 15, 15])} si va lua jucatorul
care are suma listei egala cu 210.


\subsubsection*{Exercitiul 2}

Avand ca si context urmatoarele fapte:
\begin{lstlisting}
% mina(X, Y).
mina(1, 1).
mina(1, 3).
mina(2, 2).
mina(3, 1).
\end{lstlisting}

Definiti predicatul \texttt{mine\_vecine/3} care primeste coordonatele X si Y (primii doi parametrii) si returneaza numarul de mine vecine prin al treilea parametru.

Rulati si confirmati:

\texttt{?- mine\_vecine(1, 2, R).}

\texttt{R = 3.}

\texttt{?- mine\_vecine(1, 4, R).}

\texttt{R = 1.}

\texttt{?- mine\_vecine(1, 5, R).}

\texttt{R = 0.}


\subsubsection*{Solutie:}

\begin{lstlisting}
vecin(X1,Y1, X2,Y2) :-
    between(-1, 1, Dx),
    between(-1, 1, Dy),
    \+ (Dx = 0, Dy = 0), % Sau (Dx =\= 0; Dy =\= 0)
    X2 is X1 + Dx,
    Y2 is Y1 + Dy.

mine_vecine(X, Y, R) :-
    findall((Vx,Vy), (vecin(X,Y, Vx,Vy), mina(Vx,Vy)), L),
    length(L, R).
\end{lstlisting}


\newpage

\subsection*{Exercitii propuse:}
\subsubsection*{Exercitiul 1}
Definiti predicatul \texttt{genAllBetween/3} care genereaza toate numerele dintre A si B (primii doi parametrii) si pune rezultatul ca o lista in ultimul parametru. (Hint: folositi \texttt{between/3} si \texttt{findall/3})

Rulati si confirmati:

\texttt{?-genAllBetween(4, 10, R).}

\texttt{R = [4, 5, 6, 7, 8, 9, 10]}

\texttt{?-genAllBetween(4, 4, R).}

\texttt{R = [4]}

\texttt{?-genAllBetween(5, 4, R).}

\texttt{R = []}




\subsubsection*{Exercitiul 2}
Definiti predicatul \texttt{aboveMean/2} care gaseste elementele mai mari decat media tuturor elementelor din prima lista si le pune in a doua lista. Implementati folosind una dintre predicatele \texttt{findall/3}, \texttt{bagof/3} sau \texttt{setof/3}

Rulati si confirmati:

\texttt{?-aboveMean([3, 7, 2, 8, 5, 10], R).}

\texttt{R = [7, 8, 10].}

\texttt{?-aboveMean([1, 1, 1, 1, 1, 10], R).}

\texttt{R = [10].}



\subsubsection*{Exercitiul 3}
Avand ca si context urmatoarele fapte:
\begin{lstlisting}
% nota(Student, Materie, Nota).
nota(andrei, mate1, 8).
nota(andrei, fizica, 4).
nota(andrei, programare1, 10).
nota(maria, mate1, 6).
nota(maria, fizica, 7).
nota(maria, programare1, 10).

% cerinta(MaterieDorita, MaterieCeTrebuiePromovata).
cerinta(mate2, mate1).
cerinta(programare2, programare1).
cerinta(robotica, programare1).
cerinta(robotica, fizica). % Robotica are doua cerinte
\end{lstlisting}

Definiti predicatul \texttt{meterii\_eligibile/2} care primeste numele unui student ca primul parametru si in al doilea parametru se afla lista materiilor la care poate sa se inscrie.

Rulati si confirmati:

\texttt{?- materii\_eligibile(andrei, X).}

\texttt{X = [mate2, programare2].}

\texttt{?- materii\_eligibile(maria, X).}

\texttt{X = [mate2, programare2, robotica].}
